% !TEX program = pdflatex
\documentclass[11pt,a4paper]{article}

% =========================================================================
% VIABILIDAD Y BUSINESS CASE
% Proyecto: Horus Sentinel
% Gestión de Proyectos TICS (EST315) -- Semestre 2026-II -- Team Sobrevivientes
% =========================================================================

% --- PAQUETES DE CONFIGURACIÓN ---
\usepackage[utf8]{inputenc}
\usepackage[T1]{fontenc}
\usepackage[spanish,es-tabla]{babel}
% Compatibilidad con babel-spanish reciente (flechas de TikZ y signo de porcentaje)
\AtBeginDocument{\shorthandoff{<>}\spanishplainpercent}
\usepackage{geometry}
\geometry{
    top=3.6cm,
    bottom=2.3cm,
    left=2.2cm,
    right=2.2cm,
    % PARTE MODIFICABLE: headheight debe ser mayor que la altura del logo más alto (2.56cm = 73pt).
    headheight=80pt,
    headsep=15pt,
    footskip=20pt
}

\usepackage{graphicx}
\usepackage{fancyhdr}
\usepackage[table]{xcolor}
\usepackage{titlesec}
\usepackage{booktabs}
\usepackage{tabularx}
\usepackage{array}
\usepackage{amsmath}
\usepackage{enumitem}
\usepackage{microtype}
\usepackage{amssymb}
\usepackage{hyperref}

% Logos: funciona a 2 o 3 niveles de profundidad dentro de docs/latex/
\graphicspath{{../../assets/img/}{../../../assets/img/}}

% --- PALETA INSTITUCIONAL FIEI / UNA PUNO ---
\definecolor{AzulUNA}{RGB}{10, 45, 90}
\definecolor{DoradoFIEI}{RGB}{180, 130, 20}
\definecolor{GrisOscuro}{RGB}{50, 50, 50}
\definecolor{AzulClaro}{RGB}{228, 236, 246}

% --- CONFIGURACIÓN DE LISTAS ---
\setlist{topsep=2pt, itemsep=2pt, parsep=1pt}

% --- COLUMNAS DE TABLA ---
\newcolumntype{L}{>{\raggedright\arraybackslash}X}
\newcolumntype{P}[1]{>{\raggedright\arraybackslash}p{#1}}
\newcolumntype{C}[1]{>{\centering\arraybackslash}p{#1}}
\newcolumntype{Q}[1]{>{\raggedleft\arraybackslash}p{#1}}
\newcolumntype{R}{>{\raggedleft\arraybackslash}X}
\renewcommand{\arraystretch}{1.25}

% --- FORMATO DE ENLACES ---
\hypersetup{
    colorlinks=true,
    linkcolor=AzulUNA,
    urlcolor=AzulUNA,
    citecolor=AzulUNA,
    pdftitle={Viabilidad y Business Case -- Horus Sentinel},
    pdfauthor={Team Sobrevivientes}
}

% --- FORMATO DE TÍTULOS ---
\titleformat{\section}{\color{AzulUNA}\normalfont\large\bfseries}{\thesection.}{0.7em}{}[\color{AzulUNA}\titlerule]
\titleformat{\subsection}{\color{AzulUNA}\normalfont\normalsize\bfseries}{\thesubsection.}{0.5em}{}
\titlespacing*{\section}{0pt}{9pt}{4pt}
\titlespacing*{\subsection}{0pt}{6pt}{2pt}

% --- ENCABEZADO Y PIE DE PÁGINA (fancyhdr) ---
\pagestyle{fancy}
\fancyhf{}

% ===== TAMAÑO DE LOS LOGOS (PARTE MODIFICABLE) =====
% logo_fiei.png  : height=1.45cm (lado izquierdo)
% logo_equipo.png: height=2.56cm (lado derecho; 25 % más grande porque la imagen tiene más borde blanco)
% Para cambiar el tamaño, modifique solo el valor height=... de cada \includegraphics.
% Si el logo más alto supera 2.6cm, aumente también headheight y top en la configuración de página.
\fancyhead[L]{%
    \includegraphics[height=1.45cm,keepaspectratio]{logo_fiei.png}%
}
\fancyhead[C]{%
    \scriptsize \textbf{\textcolor{AzulUNA}{UNIVERSIDAD NACIONAL DEL ALTIPLANO DE PUNO}}\\[2pt]
    \tiny \textcolor{GrisOscuro}{FACULTAD DE INGENIERÍA ESTADÍSTICA E INFORMÁTICA}\\[2pt]
    \scriptsize \textbf{\textcolor{AzulUNA}{Gestión de Proyectos TICS (EST315) \textbar\ Team Sobrevivientes}}\\[2pt]
    \footnotesize \textbf{\textcolor{DoradoFIEI}{VIABILIDAD Y BUSINESS CASE}}%
}
\fancyhead[R]{%
    \includegraphics[height=2.56cm,keepaspectratio]{logo_equipo.png}%
}

\renewcommand{\headrulewidth}{0pt}
\renewcommand{\headrule}{}

\fancyfoot[L]{\scriptsize \textcolor{GrisOscuro}{Horus Sentinel}}
\fancyfoot[C]{\small \thepage}
\fancyfoot[R]{\scriptsize \textcolor{GrisOscuro}{Puno, Perú -- Semestre 2026-II}}
\renewcommand{\footrulewidth}{0.4pt}

\begin{document}

\begin{center}
    {\Large \textbf{\textcolor{AzulUNA}{VIABILIDAD Y BUSINESS CASE}}}\\[3pt]
    {\normalsize \textbf{Horus Sentinel}}\\[4pt]
    \footnotesize \textbf{Fecha:} jueves 10/09/2026 (semana 3) \quad \textbar\quad \textbf{Versión:} 1.0 \quad \textbar\quad \textbf{Fase:} Inicio
\end{center}

\vspace{0.1cm}

\section{Datos Generales}
\noindent\begin{tabularx}{\textwidth}{@{}P{4.4cm} L@{}}
\toprule
\textbf{Proyecto:} & Horus Sentinel \\
\textbf{Responsable:} & Ruth Karina Apaza Solis -- Analista de Planificación \\
\textbf{Sponsor:} & Team Sobrevivientes, dueño y patrocinador del proyecto \\
\textbf{Zona piloto:} & Puno (cercado) \\
\textbf{Tipo de estimación (PM4DEV):} & Aproximada, con un rango de $-25\,\%$ a $+75\,\%$ \\
\textbf{Horizonte de evaluación:} & Semestre 2026-II (inversión) y años 1 a 3 (operación); tasa de descuento de 10\,\% \\
\bottomrule
\end{tabularx}

\section{Matriz de Viabilidad}
{\small
\noindent\begin{tabularx}{\textwidth}{@{}P{2.3cm} L P{2.6cm}@{}}
\toprule
\rowcolor{AzulClaro}
\textbf{Dimensión} & \textbf{Análisis} & \textbf{Dictamen} \\
\midrule
\textbf{Técnica} & El mapa de calor (KDE) y la ruta segura ($A^*$ con costo $\alpha \cdot$ distancia $+ \beta \cdot$ riesgo) se construyen con mapas abiertos de Puno. El equipo no tiene experiencia en análisis geoespacial y el cálculo requiere un servidor con más memoria. & Viable con condiciones: capacitación geoespacial y servidor de mayor capacidad \\
\addlinespace[2pt]
\textbf{Económica} & Costo valorizado del semestre: S/ 11,829.93, con S/ 3,229.30 de gasto en efectivo (S/ 807.33 por integrante). El convenio municipal de S/ 500 al mes no alcanza a cubrir la operación anual (S/ 6,182.75). & Viable con condiciones: firma del convenio con la municipalidad \\
\addlinespace[2pt]
\textbf{Operativa} & El prototipo usa datos simulados. El MVP depende de reportes ciudadanos, que requieren 3 h semanales de moderación y una campaña de difusión en el cercado. & Viable con condiciones \\
\addlinespace[2pt]
\textbf{Legal} & Solo incidentes agregados por zona y horario, sin identificar personas (Ley N.\textsuperscript{o} 29733). El convenio con la Municipalidad Provincial de Puno requiere un acuerdo formal. & Viable con condiciones \\
\addlinespace[2pt]
\textbf{Temporal} & Prototipo en la semana 8 y MVP en la semana 16, con 8 h semanales por integrante. El cronograma no deja holgura y comparte al equipo con CasaScore. & Viable con condiciones: riesgo alto de retraso \\
\bottomrule
\end{tabularx}
}

\section{Costos del Proyecto (semestre 2026-II)}
Mismos parámetros que CasaScore: tarifas por hora de mercado con EsSalud, 8 h semanales por integrante durante 16 semanas, 6 jornadas nocturnas y depreciación de equipos al 25\,\% anual.

\vspace{0.1cm}
{\small
\noindent\begin{tabularx}{\textwidth}{@{}P{3.6cm} L Q{2.0cm}@{}}
\toprule
\rowcolor{AzulClaro}
\textbf{Rubro} & \textbf{Detalle} & \textbf{S/} \\
\midrule
Personal regular & 512 h del equipo & 6,176.00 \\
Horas extra & 6 jornadas nocturnas $\times$ 4 h $\times$ 4 integrantes & 1,505.40 \\
Equipos & Depreciación de 4 laptops y 1 celular de pruebas & 919.23 \\
Servidor geoespacial & Servidor virtual de 4 núcleos y 8 GB por 4 meses (US\$ 40 al mes) & 549.76 \\
Servicios & Electricidad, internet del hogar y datos móviles & 540.54 \\
Alimentación & Dietas nocturnas, refrigerios y almuerzos en 8 días de campo & 1,240.00 \\
Productos de salud & Gotas para los ojos, multivitamínico, analgésico y bloqueador solar & 314.00 \\
Transporte & Taxi nocturno y pasajes de campo en el cercado & 240.00 \\
Documentación y difusión & Impresiones, 40 afiches para captar reportes e incentivos para 30 participantes & 345.00 \\
\midrule
\textbf{Estimación de costos} & Rango PM4DEV: S/ 8,872.45 a S/ 20,702.38 & \textbf{11,829.93} \\
(+) Contingencia & Valor esperado de los riesgos identificados & 591.50 \\
\textbf{Línea base de costos} & & \textbf{12,421.43} \\
(+) Reserva de gestión 10\,\% & Imprevistos & 1,242.14 \\
\textbf{Presupuesto total} & & \textbf{13,663.57} \\
\bottomrule
\end{tabularx}
}

\section{Ingresos y Costos de Operación (años 1 a 3)}
{\small
\noindent\begin{tabularx}{\textwidth}{@{}P{3.4cm} L Q{1.9cm}@{}}
\toprule
\rowcolor{AzulClaro}
\textbf{Concepto} & \textbf{Supuesto} & \textbf{S/ por año} \\
\midrule
Convenio municipal & Municipalidad Provincial de Puno: S/ 500 al mes por mantenimiento y reportes & 6,000.00 \\
\midrule
Servidor geoespacial & US\$ 40 al mes & 1,649.28 \\
Mantenimiento & 4 h/semana $\times$ S/ 12.49 & 2,597.92 \\
Moderación de reportes & 3 h/semana $\times$ S/ 10.22 & 1,594.32 \\
Dominio y difusión & Dominio web y campaña anual & 341.23 \\
\textbf{Costos de operación} & & \textbf{6,182.75} \\
\bottomrule
\end{tabularx}
}

\section{Evaluación del Business Case}

\subsection{Evaluación financiera}
{\small
\noindent\begin{tabularx}{\textwidth}{@{}C{1.2cm} R R R@{}}
\toprule
\rowcolor{AzulClaro}
\textbf{Año} & \textbf{Flujo neto} & \textbf{Descontado (10\,\%)} & \textbf{Acumulado} \\
\midrule
0 & $-$11,829.93 & $-$11,829.93 & $-$11,829.93 \\
1 & $-$182.75 & $-$166.14 & $-$12,012.68 \\
2 & $-$182.75 & $-$151.03 & $-$12,195.43 \\
3 & $-$182.75 & $-$137.30 & $-$12,378.18 \\
\midrule
\multicolumn{2}{@{}l}{\textbf{VAN financiero}} & \textbf{$-$12,284.40} & \\
\bottomrule
\end{tabularx}
}

\subsection{Evaluación económica (ingresos más beneficio para los ciudadanos)}
Beneficio social: 1,000 usuarios activos, de los cuales el 3\,\% evita un robo con un valor promedio de S/ 400 al elegir rutas más seguras: S/ 12,000 al año.

\vspace{0.1cm}
{\small
\noindent\begin{tabularx}{\textwidth}{@{}P{4.0cm} L Q{3.1cm}@{}}
\toprule
\rowcolor{AzulClaro}
\textbf{Indicador} & \textbf{Cálculo} & \textbf{Resultado} \\
\midrule
Flujo económico anual & $-$182.75 + 12,000.00 & S/ 11,817.25 \\
VAN económico & $-$11,829.93 + 10,742.95 + 9,766.32 + 8,878.47 & S/ +17,557.81 \\
ROI (3 años) & (35,451.75 $-$ 11,829.93) / 11,829.93 & 199.7\,\% \\
Payback & 11,829.93 / 11,817.25 & 1.00 año \\
Relación beneficio/costo & 29,387.74 / 11,829.93 & 2.48 \\
TIR & Tasa que hace VAN = 0 & 83.8\,\% \\
\bottomrule
\end{tabularx}
}

\section{Firmas}
\noindent En señal de conformidad con el presente análisis, los integrantes del equipo firman a continuación:

\vspace{1.5cm}
\begin{center}
\begin{tabular}{c c c}
    \parbox{6.5cm}{\centering \rule{6.0cm}{0.6pt}\\[5pt] \textbf{Clyde Neil Paricahua Pari}\\[2pt] {\footnotesize Project Manager}}
    & \hspace{0.8cm} &
    \parbox{6.5cm}{\centering \rule{6.0cm}{0.6pt}\\[5pt] \textbf{Ruth Karina Apaza Solis}\\[2pt] {\footnotesize Analista de Planificación}}
    \\[2.0cm]
    \parbox{6.5cm}{\centering \rule{6.0cm}{0.6pt}\\[5pt] \textbf{Eliseo Tarqui Ajahuana}\\[2pt] {\footnotesize Especialista Técnico / Arquitectura}}
    & \hspace{0.8cm} &
    \parbox{6.5cm}{\centering \rule{6.0cm}{0.6pt}\\[5pt] \textbf{Viviana Mary Quisbert Quispe}\\[2pt] {\footnotesize Gestora de Comunicación, Riesgos y Documentación}}
\end{tabular}
\end{center}
\end{document}
